%%%%%%%%%%%%%%%
%
% $Autor: Leon Ritscher %
% $Datum: 2025-10-07 %
% $Pfad: General/Batterie.tex
% $Version: 1785 $
%
%
%%%%%%%%%%%%%%%

\chapter{Batterie}
\label{ch:Batterie}

Um das Mikrocontroller-System unabhängig vom Computer zu machen kann eine Batterie genutzt werden.

\section{9 V-Blockbatterie}

Die 9 V Blockbatterie, welche man in der Abbildung \ref{fig:9V_Blockbatterie} erkennt, dient als primäre Energiequelle für das gesamte Messsystem. Sie ermöglicht den mobilen Betrieb aller Komponenten wie Mikrocontroller, Sensoren, OLED-Display, Real Time Clock und SD-Kartenmodul. Über den zweipoligen Wippschalter kann die Versorgung manuell aktiviert oder deaktiviert werden. Typischerweise handelt es sich bei den 9 V Blockbatterien um Alkine oder Lithium-Batterien, welche im Gegensatz zu Akkus nicht wiederaufladbar sind.

\begin{figure}[h]
	\begin{center}
		\includegraphics[width=3in]{Battery/Battery.png}
		\caption{Druckknopfanschluss für 9V Blockbatterie}
		\label{fig:9V_Blockbatterie}
	\end{center}
\end{figure} 

\section{Testen der 9 V Blockbatterie}

Um die Batteriespannung der Batterie messen zu können, kann das Programm \ref{lst:TestBattery} verwendet werden. Die Batteriespannung wird dabei über einen Spannungsteiler und den analogen Pin gemessen. Die Batterie wird während des Tests an den Batterieclip angeschlossen und der gemessene Wert bei dem seriellen Monitor ausgelesen.\\
Für den Test werden folgende Komponenten benötigt:

\begin{itemize}
	\item Arduino Nano 33 BLE Sense Lite
	\item USB-A auf Micro-USB Verbindungskabel
	\item 9 V Blockbatterie
	\item Batterieclip (9 V, mit Anschlussleitungen)
	\item 2-poliger Wippschalter (Ein/Aus)
	\item Spannungsteiler aus zwei Widerständen (z.\,B.\ 100~k$\Omega$ und 33~k$\Omega$)

\end{itemize}

\subsection{Aufbau}

Die 9 V Blockbatterie wird über einen Wippschalter mit dem System verbunden. Zur Messung der Batteriespannung ist ein Spannungsteiler vorgeschaltet, der die Spannung auf einen für den analogen Eingang (Pin A7) des Mikrocontrollers ungefährlichen Bereich reduziert. Die Spannung an Pin A7 überschreitet dabei nicht 3,3 V.

\subsection{Durchführung}

Der Arduino Nano 33 BLE Sense Lite wird über USB mit dem Computer verbunden. Anschließend wird der folgende Sketch auf den Mikrocontroller geladen und der serielle Monitor in der Arduino~IDE geöffnet. Die Batterie wird über den Schalter aktiviert. Der aktuelle Spannungswert wird nun in 5-Sekunden-Intervallen im seriellen Monitor angezeigt.

\subsection{Ergebnisse}

\begin{code}[h]
	\captionof{code}{Testen der 9 V Blockbatterie} \label{lst:TestBattery}
	\pythonexternal[language=c++]{../../Code/Arduino/Battery/Test/Test.ino}
\end{code}

Zu Beginn des Tests zeigt der serielle Monitor die Spannung 0 V an. Nach dem Anschließen der Batterie an den Spannungssensor wird die Spannung 9,6 V angezeigt.  Abbildung~\ref{fig:BildSpannungTest} zeigt den gemessenen Spannungsverlauf. Die angezeigte Spannung liegt in dem erwarteten Bereich einer 9 V Batterie.

\begin{figure}[h]
	\begin{center}
		\includegraphics[width=\textwidth]{Battery/BatteryTest.png}
		\captionof{figure}{Output am seriellen Monitor der Spannungsüberwachung}
		\label{fig:BildSpannungTest}
	\end{center}
\end{figure}

% ---------------------------------------------------------------------------------------------------------------------------- %

\section{9 V Blockbatterie Druckknopfanschluss}
Der Druckknopfanschluss in Abbildung \ref{fig:9V_Druckknopfanschluss} kann vertikal an eine 9 V Blockbatterie angeschlossen werden. Er verfügt über eine Folienausführung und einen T-Ausführungsanschluss, was eine einfache und sichere Verbindung ermöglicht. Die dazugehörigen Anschlussdrähte haben eine Länge von 150 mm.

\begin{figure}[h]
    \begin{center}
        \includegraphics[width=3in]{Battery/clip}
        \caption{Druckknopfanschluss für 9V Blockbatterie}
        \label{fig:9V_Druckknopfanschluss}
    \end{center}
\end{figure} 

\subsubsection{Technische Daten}
\begin{description}
	\item[Passend für Batteriegröße:] 9 V Block
	\item[Abmessungen (Länge x Breite x Höhe):] 26mm x 13mm x 8mm
	\item[Ausführung:] T-Ausführung
	\item[Kabellänge:] 150mm
\end{description}